\documentclass[10pt,a4paper]{amsart}
\usepackage[margin=19mm]{geometry}
\usepackage{amsmath,amssymb,mathtools}
\usepackage[scr=rsfs,cal=euler]{mathalfa}
\usepackage{xfrac}
\usepackage[unicode,hidelinks,bookmarks=false]{hyperref}
\newcommand{\ie}{i.e.\ }
\newcommand{\cf}{cf.\ }
\newcommand{\resp}{respectively\ }
\newcommand{\Subsection}{\subsection}
\providecommand{\nobreakdash}{\nobreak\hbox{-}}
\setlength{\emergencystretch}{2em}
\multlinegap0pt
\usepackage[shortlabels]{enumitem}
\usepackage{amssymb}
\usepackage{xcolor}
\usepackage{mathtools}
\usepackage{booktabs}
\usepackage{dsfont}
\newtheorem{thm}{Theorem}
\newcommand{\FF}{\mathcal{F}}
\newcommand{\ff}{\mathcal{F}}
\newcommand{\m}{\mathcal}
\title{Frankl's diversity theorem for permutations}
\author{Eduard Inozemtsev, Andrey Kupavskii}
\date{Source: arXiv:2602.11796v1 (CC BY 4.0)}
\begin{document}
\maketitle

\noindent\emph{Source and licence.} Frankl's diversity theorem for permutations. Version arXiv:2602.11796v1 (CC BY 4.0), distributed by arXiv under the Creative Commons Attribution 4.0 International licence, http://creativecommons.org/licenses/by/4.0/, as recorded in the arXiv licence metadata for this exact version and shown on the abstract page https://arxiv.org/abs/2602.11796v1. Full licence notice: \texttt{licenses/arxiv-2602.11796-CC-BY-4.0.txt}.\par\smallskip

\noindent\textit{Editorial scope.} Counted unit: the article's thm thmapprox (source0.tex lines 462--470). The article's complete original proof of it is delivered with it (source0.tex lines 471--487, one \texttt{proof} environment of 17 lines). No mathematics is written, repaired or re-derived: the statement and the proof are the article's own lines, in the article's own notation, and no retained line is substituted or re-flowed.  No further source-local context is needed: every cross-reference of every retained line resolves inside the two delivered spans.  Nothing was omitted from any retained span, so the package carries no omission record.  No retained line contains a citation, so the wrapper carries no bibliography and no bibliography entry.  No retained line contains a figure environment, an image-inclusion command or drawing code, so no image is delivered and the package declares no asset dependency.  Editorial formatting only, in addition: an AMS \texttt{amsart} preamble that restates the article's own declarations and macros used by the retained lines, namely the environments \texttt{thm} on separate counters, together with the standard packages those lines need.  The retained environments print in this document as Theorem 1 in order of appearance, whereas the article prints the same items with the numbering of its own full document.  The complete native source of the article as distributed in the arXiv e-print for this exact version is bundled unchanged as \texttt{sources/194491/source0.tex}.
\begin{thm}\label{thmapprox}
    Consider a family $\FF_1$  of permutations on $[2,n]$ and a family $\ff_2$ of partial permutations on $[2,t]$, such that $\FF_1$ and $\FF_2$ are cross-intersecting. 
    Let $\tau_1 := 1.01, \tau_2:=\frac{k}{\log^2{n}}$ and let $q_1:=10\log{n}, q_2:=\lceil1.1(t-k)\rceil.$ Then there exist families $\m{S}_1, \m S_2$ of partial permutations on  $[2,n]$ and $[2,t]$ of sizes at most $q_1,q_2$, respectively, and families $\FF_i' \subset \FF_i$ for $i\in\{1,2\}$ such that the following holds.
    \begin{enumerate}[label=(\roman*)]
        \item We have $\FF_i \setminus \FF_i'\subset \Sigma_{n_i}[\m S_i]$;
        \item for any $B \in \m{S}_i$ there is a family $\FF_{i, B} \subset \FF_i$ such that $\FF_{i,B}(B)$ is $(\Sigma_{n_i}(B), \tau_i)$-homogeneous;
        \item $|\FF_i'| \leq \tau_i^{-q_i-1}n_i!$.
    \end{enumerate}
\end{thm}

\begin{proof} The first guiding observation to make is that $\FF$ cannot be $\tau$-homogeneous.  
Indeed, the same application of Spread Lemma as we will use later imply that there are $A_1 \in \FF_1, A_2 \in \FF_2$ such that $A_1 \cap A_2=\emptyset$. We avoid calculations here, since we do them below in a more general scenario.  

\begin{enumerate}
    \item Find an inclusion-maximal $S_i^{j}$ that  $|\FF_{i}^j(S_i^{j})| \ge  \tau^{|S_i^{j}|}_i \frac{|\Sigma_{n_i}(S)|}{|\Sigma_{n_i}|}|\FF_i^{j}|$.
    \item If $|S_i^{j}|> q_i$ or $\FF_{i}^j = \emptyset$ then stop. Otherwise, put $\FF_{i}^{j+1}:=\FF_{i}^j\setminus \FF_{i}^{j}[S_i^{j}]$. 
\end{enumerate}
Let us show that $\FF_{i}^j(S_i^{j})$ is $(\Sigma_{n_i}(B), \tau_i)$-homogeneous. Indeed, for any set $X$ disjoint with $S_i^j$ we have
$$|\FF_{i}^j(S_i^{j}\cup X)|< \tau_i^{|S_i^{j}|+|X|} \frac{|\Sigma_{n_i}(S_i^j\cup X)|}{|\Sigma_{n_i}(X)|} |\ff_i^j|\le \tau_i^{|X|} \frac{|\Sigma_{n_i}(S_i^j)|}{|\Sigma_{n_i}|} |\FF_{i}^j (S_i^{j})|,$$
where the first inequality uses the maximality of $S_i^j.$

Let $N_i$ be the step of the procedure for $\FF_i$ at which we stop. The families $\m{S}_i$ are defined as follows: $\m{S}_i:=\{S_i^{1},\ldots, S_i^{N_i-1}\}$. Clearly, $|S_i^{j}|\le q_i$ for each $j\in [N_i-1]$. The families $\FF_{B}$ promised in (ii) are defined to be $\FF_{i}^j[S_i^j]$ for $B=S_i^j$. Next, note that if $\FF_i^{N_i}$ is non-empty, then 
$$
|\FF_{i}^{N_i}|\le \tau_i^{-|S_i^{N_i}|} \frac{|\Sigma_{n_i}|}{|\Sigma_{n_i}(S)|}|\FF_i(S_i^{N_i})|\le \tau_i^{-|S_i^{N_i}|}|\Sigma_{n_i}|.
$$
We put $\FF_i':=\FF_i^{N_i}$. Since either $|S_i^{N_i}|>q_i$ or $\FF_i' = \emptyset $, we have $|\FF_i'| \leq \tau^{-q_i-1}_i \cdot n_i!$.
\end{proof}

\end{document}
